\documentclass[10pt,a4paper]{amsart}
\usepackage[margin=19mm]{geometry}
\usepackage{amsmath,amssymb,mathtools}
\usepackage[scr=rsfs,cal=euler]{mathalfa}
\usepackage{xfrac}
\usepackage[unicode,hidelinks,bookmarks=false]{hyperref}
\newcommand{\ie}{i.e.\ }
\newcommand{\cf}{cf.\ }
\newcommand{\resp}{respectively\ }
\newcommand{\Subsection}{\subsection}
\providecommand{\nobreakdash}{\nobreak\hbox{-}}
\setlength{\emergencystretch}{2em}
\multlinegap0pt
\usepackage{booktabs}
\usepackage{nicefrac}
\usepackage{xcolor}
\usepackage[shortlabels]{enumitem}
\usepackage{float}
\usepackage{algorithm}\usepackage{algpseudocode}
\usepackage{bbm}
\newtheorem{proposition}{Proposition}
\title{Causal LLM Routing: End-to-End Regret Minimization from Observational Data}
\author{Asterios Tsiourvas, Wei Sun, Georgia Perakis}
\date{Source: arXiv:2505.16037v2 (CC BY 4.0)}
\begin{document}
\maketitle

\noindent\emph{Source and licence.} Causal LLM Routing: End-to-End Regret Minimization from Observational Data. Version arXiv:2505.16037v2 (CC BY 4.0), distributed by arXiv under the Creative Commons Attribution 4.0 International licence, http://creativecommons.org/licenses/by/4.0/, as recorded in the arXiv licence metadata for this exact version and shown on the abstract page https://arxiv.org/abs/2505.16037v2. Full licence notice: \texttt{licenses/arxiv-2505.16037-CC-BY-4.0.txt}.\par\smallskip

\noindent\textit{Editorial scope.} Counted unit: the article's proposition proposition (source0.tex lines 815--828). The article's complete original proof of it is delivered with it (source0.tex lines 830--875, one \texttt{proof} environment of 46 lines). No mathematics is written, repaired or re-derived: the statement and the proof are the article's own lines, in the article's own notation, and no retained line is substituted or re-flowed.  No further source-local context is needed: every cross-reference of every retained line resolves inside the two delivered spans.  Nothing was omitted from any retained span, so the package carries no omission record.  No retained line contains a citation, so the wrapper carries no bibliography and no bibliography entry.  No retained line contains a figure environment, an image-inclusion command or drawing code, so no image is delivered and the package declares no asset dependency.  Editorial formatting only, in addition: an AMS \texttt{amsart} preamble that restates the article's own declarations and macros used by the retained lines, namely the environments \texttt{proposition} on separate counters, together with the standard packages those lines need.  The retained environments print in this document as Proposition 1 in order of appearance, whereas the article prints the same items with the numbering of its own full document.  The complete native source of the article as distributed in the arXiv e-print for this exact version is bundled unchanged as \texttt{sources/178551/source0.tex}.
\begin{proposition}
Let $f: \mathcal{X} \to \mathbb{R}^{|\mathcal{T}|}$ be a neural network whose output is passed through a softmax layer with fixed temperature $\tau > 0$, and define $t^*_i := \arg\max_{t \in \mathcal{T}} \hat{Y}_{x_i}(t)$. Then, optimizing  the following objective using gradient descent 
\begin{align}
    \min_f \frac{1}{n} \sum_{i=1}^n \left( \hat{Y}_{x_i}(t^*_i) - \sum_{t=1}^{|\mathcal{T}|} \hat{Y}_{x_i}(t) \cdot \mathrm{softmax}(f(x_i))_t \right)
\end{align}
leads the model $f$ to place all probability mass on the optimal treatment $t^*_i$. That is, at convergence,
\begin{align}
    \mathrm{softmax}(f(x_i))_t \to 
    \begin{cases}
        1 & \text{if } t = t^*_i, \\
        0 & \text{otherwise}.
    \end{cases}
\end{align}
\end{proposition}

\begin{proof}
Let $\hat{Y}_{x_i} \in \mathbb{R}^{|\mathcal{T}|}$ denote the vector of estimated potential outcomes for input $x_i$, and let $f(x_i) \in \mathbb{R}^{|\mathcal{T}|}$ be the output of the neural network before the softmax layer. The objective for a single instance $x_i$ can be written as minimizing the regret surrogate:
\begin{align}
    \hat{Y}_{x_i}(t^*_i) - \sum_{t=1}^{|\mathcal{T}|} \hat{Y}_{x_i}(t) \cdot \mathrm{softmax}(f(x_i))_t.
\end{align}
This is equivalent to maximizing the inner product:
\begin{align}
    \langle \hat{Y}_{x_i}, \mathrm{softmax}(f(x_i)) \rangle.
\end{align}

Let us denote $p := \mathrm{softmax}(f(x_i)) \in \Delta^{|\mathcal{T}| - 1}$, the probability simplex. We now show that the inner product $\langle \hat{Y}_{x_i}, p \rangle$ increases at each gradient step. Since $p = \mathrm{softmax}(f(x_i))$, we can compute the gradient of the objective with respect to $f(x_i)$ as:
\begin{align}
    \nabla_{f(x_i)} \langle \hat{Y}_{x_i}, \mathrm{softmax}(f(x_i)) \rangle
    = J_{\mathrm{softmax}}(f(x_i))^\top \hat{Y}_{x_i},
\end{align}
where $J_{\mathrm{softmax}}(f(x_i))$ is the Jacobian of the softmax function, given by:
\begin{align}
    J_{\mathrm{softmax}}(f(x_i))_{t,s} = \frac{\partial \, \mathrm{softmax}(f(x_i))_t}{\partial f(x_i)_s}
    = \mathrm{softmax}(f(x_i))_t \left(\delta_{t,s} - \mathrm{softmax}(f(x_i))_s \right).
\end{align}

This gradient direction corresponds to increasing the logit value of actions with higher $\hat{Y}_{x_i}(t)$ and decreasing those with lower values, pushing the softmax distribution toward the mode of $\hat{Y}_{x_i}$. In other words, the gradient ascent step increases the inner product at each iteration $k$:
\begin{align}
    \langle \hat{Y}_{x_i}, \mathrm{softmax}(f(x_i)) \rangle^{(k+1)} 
    > \langle \hat{Y}_{x_i}, \mathrm{softmax}(f(x_i)) \rangle^{(k)}.
\end{align}

Since $\hat{Y}_{x_i}$ is fixed and the softmax is smooth and bounded, this sequence is monotonically increasing and converges to the maximum possible value:
\begin{align}
    \langle \hat{Y}_{x_i}, \mathrm{softmax}(f(x_i)) \rangle \to \max_t \hat{Y}_{x_i}(t) = \hat{Y}_{x_i}(t^*_i),
\end{align}
which implies:
\begin{align}
    \mathrm{softmax}(f(x_i))_t \to 
    \begin{cases}
        1 & \text{if } t = t^*_i, \\
        0 & \text{otherwise}.
    \end{cases}
\end{align}

Thus, the regret surrogate converges to zero:
\begin{align}
    \hat{Y}_{x_i}(t^*_i) - \langle \hat{Y}_{x_i}, \mathrm{softmax}(f(x_i)) \rangle \to 0,
\end{align}
and the learned policy selects the treatment maximizing the estimated outcome.
\end{proof}

\end{document}
